\begin{annexitem}{annex:p07-circle-integer-presentation}{thm:p07-circle-integer-presentation}{Natural-number codes and canonical integer quotient}
\begin{proof}
Here are the natural-number facts needed by the integer construction, rather
than an assumed set structure on an unspecified integer type. Define
\(\mathrm{EqN}(0,0)=1\),
\(\mathrm{EqN}(0,\mathsf Sb)=0_{\mathrm{type}}\),
\(\mathrm{EqN}(\mathsf Sa,0)=0_{\mathrm{type}}\), and
\(\mathrm{EqN}(\mathsf Sa,\mathsf Sb)=\mathrm{EqN}(a,b)\).
The diagonal element \(e_a:\mathrm{EqN}(a,a)\) is defined by induction,
with \(e_0=\star\) and \(e_{\mathsf Sa}=e_a\).
Encode \(p:a=b\) by transport of \(e_a\) in the second coordinate.
Decode by double induction: at \(0,0\) eliminate the unit and return
reflexivity; mixed constructors eliminate the empty type; at
\(\mathsf Sa,\mathsf Sb\) apply \(\operatorname{ap}_{\mathsf S}\) to the
recursive decoder.

Decoding an encoded path is the path itself. Indeed, based \(J\) reduces
this assertion to decoding \(e_a\) as reflexivity. Induction on \(a\)
gives it at zero by unit computation; at \(\mathsf Sa\) use the induction
path and \(\operatorname{ap}_{\mathsf S}(\mathsf{refl})
\equiv_{\mathrm j}\mathsf{refl}\).
Each \(\mathrm{EqN}(a,b)\) is a proposition: use double induction, the unit
and empty proposition proofs, and the recursive clause. Hence the other
composite, encode after decode, equals the original code without assuming
anything about identity proofs in \(\mathbb N\). The two maps form an
equivalence by adjointification. Thus each \(a=b\) is a proposition,
so \(\mathbb N\) is a set. The same code gives \(0\ne\mathsf Sa\)
and injectivity of \(\mathsf S\): mixed codes are empty, and equal
successors decode an element of \(\mathrm{EqN}(a,b)\).
This expands the natural-number calculation of
\cite[\S2.13]{HoTT2013PublicReleaseFoundations}.

For clarity, take natural addition recursively in its first argument:
\(0+b=b\), \(\mathsf Sa+b=\mathsf S(a+b)\).
Induction on \(a\) proves \(a+0=a\) and
\(a+\mathsf Sb=\mathsf S(a+b)\): their zero cases are reflexivity
and their successor cases are \(\operatorname{ap}_{\mathsf S}\) of the
induction paths. Associativity follows by induction on the first argument,
with the same successor step. Commutativity follows by induction on the
first argument using the just-proved second-argument successor law and
right-zero law. Left cancellation follows by induction on the cancelled
argument, cancelling \(\mathsf S\) in the successor case. Right
cancellation follows by commutativity. These facts do not invoke integers.

For \(C=1+\mathbb N+\mathbb N\), use a code which is unit on the zero
tag against itself, \(\mathrm{EqN}(a,b)\) on equal positive or equal
negative tags, and empty on distinct tags. Tag induction defines the
diagonal code, encoder and decoder; decoding the diagonal is reflexivity
by the preceding natural calculation. The codes are propositions,
so encode after decode follows as above. This proves that \(C\) is a set
and that its three tags are disjoint. The pair type \(\mathbb N^2\)
is a set by the Sigma-path description. Each fiber \(r(x)=x\) is
therefore a proposition; consequently \(\mathbb Z\) is a set by the
same Sigma-path description and proposition closure.

Tag cases give a path \(g(j(c))=c\), with judgmental computation
on each constructor. Thus \(r(r(x))=r(x)\).
They also give \(GE(c)=c\). For \(z=(x,h):\mathbb Z\), the first
coordinates of \(EG(z)\) and \(z\) are related by \(h:r(x)=x\).
After transport along \(h\), their second coordinates lie in the same
proposition \(r(x)=x\), so they are equal. The Sigma-path constructor
gives
\(\nu_z:EG(z)=z\).
Together these paths give an equivalence \(E:C\simeq\mathbb Z\)
using the already proved adjointification theorem. In particular signed
induction at an arbitrary \(z\) is tag induction at \(G(z)\), followed by
transport along \(\nu_z\). It does not require proof-irrelevance in the
universe or in circle identity types.

We verify the advertised relation, including unequal signs. Removing
one successor from both entries of the first pair preserves
\(a+d=b+c\) by successor cancellation. Removing one successor from
both entries of the second pair preserves it by the second-argument
successor law and cancellation. Each removal also preserves the
corresponding \(g\). Induction on the two pairs reduces both to their
canonical pairs. The nine resulting tag comparisons are as follows:
zero against zero has equation \(0=0\); positive against positive has
\(\mathsf Sa=\mathsf Sc\), hence \(a=c\); negative against negative
has \(\mathsf Sd=\mathsf Sb\), hence \(b=d\).
Zero against positive, positive against zero, zero against negative and
negative against zero all assert that zero equals a successor.
Positive against negative asserts
\(\mathsf Sa+\mathsf Sd=0\), and negative against positive asserts
\(0=\mathsf Sb+\mathsf Sc\); both are impossible by the defining
addition equations and zero/successor disjointness.
In each equal-tag case, the integer code above identifies precisely
the displayed equality of indices; in each other case both sides of
the proposed equivalence are empty. This proves both directions of
\(g(a,b)=g(c,d)\) iff \(a+d=b+c\). In particular \(\sim\) is an
equivalence relation, since it is equality after \(g\).

Finally let \(B\) be a set. A function \(f:\mathbb Z\to B\)
restricts to \(h=fq\), which respects \(\sim\), by the relation just
proved. Conversely, from a relation-respecting \(h:\mathbb N^2\to B\)
define \(\bar h(x,h_x)=h(x)\).
For every \(x\), \(r(x)\sim x\), because \(gr(x)=gx\);
thus \(\bar h(q(x))=h(r(x))=h(x)\).
For \(z=(x,h_x)\), the Sigma path used above also gives \(q(x)=z\).
Consequently extending \(fq\) returns \(f\) pointwise. Funext gives
the two function equalities. The extra respect-witnesses are equal
because they are functions into identity types of the set \(B\).
This is the quotient universal property into sets, not an unsupported
identification of a signed datatype with the source quotient.
\end{proof}
\end{annexitem}

\begin{annexitem}{annex:p07-circle-integer-steps}{thm:p07-circle-integer-steps}{Successor equivalence and descended addition}
\begin{proof}
The tables supply inverse maps on \(C\) by exhaustive cases.
For \(ts\), zero goes to \(\mathsf{pos}(0)\) and back to zero;
\(\mathsf{pos}(k)\) goes to \(\mathsf{pos}(\mathsf Sk)\) and back;
\(\mathsf{neg}(0)\) goes to zero and back; and
\(\mathsf{neg}(\mathsf Sk)\) goes to \(\mathsf{neg}(k)\) and back.
For \(st\), zero goes to \(\mathsf{neg}(0)\) and back;
\(\mathsf{neg}(k)\) goes to \(\mathsf{neg}(\mathsf Sk)\) and back;
\(\mathsf{pos}(0)\) goes to zero and back; and
\(\mathsf{pos}(\mathsf Sk)\) goes to \(\mathsf{pos}(k)\) and back.
These tag computations are judgmental. Using \(GE=\mathrm{id}\)
and \(\nu:EG\sim\mathrm{id}\) gives
\(\mathsf{pred}(\mathsf{succ}(z))=z\) and
\(\mathsf{succ}(\mathsf{pred}(z))=z\) for arbitrary \(z\).
Adjointification, not a claim that every quasi-inverse automatically has
a judgmental triangle, produces \(e_{\mathsf{succ}}\).

For the pair operation \(F((a,b),(c,d))=q(a+c,b+d)\), suppose
\(a+b'=b+a'\) and \(c+d'=d+c'\). Addition associativity and
commutativity identify
\[
(a+c)+(b'+d')=(a+b')+(c+d')
 =(b+a')+(d+c')=(b+d)+(a'+c').
\]
The relation criterion therefore gives equal outputs of \(F\).
Descent twice along the universal property in the preceding proof
defines addition on \(\mathbb Z\); function types into the set
\(\mathbb Z\) are sets by funext. Its displayed pair comparison follows
from the two descent homotopies.

The identities
\[
\mathsf{succ}(q(a,b))=q(\mathsf Sa,b),\qquad
\mathsf{pred}(q(a,b))=q(a,\mathsf Sb)
\]
are checked before using the additive recurrences. For successor,
cancel common successors in \((a,b)\). On canonical pairs, zero and
positive pairs are the first two clauses of \(s\); \((0,1)\) yields
\(q(1,1)=0\); and \((0,\mathsf S(\mathsf Sk))\) yields
\(q(1,\mathsf S(\mathsf Sk))=q(0,\mathsf Sk)\), the last clause.
Cancellation commutes with adding the extra first successor:
\(g(\mathsf S(\mathsf Sa),\mathsf Sb)=g(\mathsf Sa,b)\).
This is the recursive step, so the canonical check proves the equation
for all pairs. For predecessor the canonical cases are zero and
negative pairs, \((1,0)\) giving \(q(1,1)=0\), and
\((\mathsf S(\mathsf Sk),0)\) giving
\(q(\mathsf S(\mathsf Sk),1)=q(\mathsf Sk,0)\).
Its recursive step is
\(g(\mathsf Sa,\mathsf S(\mathsf Sb))=g(a,\mathsf Sb)\).
Thus no sign boundary is omitted.

Every \(z=(x,h)\) is equal to \(q(x)\), so checking an equation of
integers on pair representatives suffices by substitution.
Right and left zero follow from natural right and left zero.
For the successor recurrence compute on \(m=q(a,b),n=q(c,d)\):
\[
m+\mathsf{succ}(n)
 =q(a+\mathsf Sc,b+d)
 =q(\mathsf S(a+c),b+d)
 =\mathsf{succ}(m+n).
\]
For predecessor compute
\[
m+\mathsf{pred}(n)
 =q(a+c,b+\mathsf Sd)
 =q(a+c,\mathsf S(b+d))
 =\mathsf{pred}(m+n).
\]
These are composites of propositional comparison paths.
\end{proof}
\end{annexitem}

\begin{annexitem}{annex:p07-circle-decoder}{thm:p07-circle-decoder}{The actual dependent loop datum of decoding}
\begin{proof}
Put \(D(x)=\mathrm{code}(x)\to(\mathsf{base}=x)\) and
\(d(n)=\mathsf{loop}^n:D(\mathsf{base})\).
For \(n:\mathbb Z\), function transport from
Theorem~\ref{thm:p07-hott-function-paths} gives
\[
T_D(\mathsf{loop},d)(n)
 =T_{x\mapsto(\mathsf{base}=x)}
   \bigl(\mathsf{loop},
          d(T_{\mathrm{code}}(\mathsf{loop}^{-1},n))\bigr).
\]
Transport in a based-path family is right concatenation, so the right
side equals
\(d(T_{\mathrm{code}}(\mathsf{loop}^{-1},n))\cdot\mathsf{loop}\).
Inverse monodromy and application of \(d\) replace its argument by
\(\mathsf{pred}(n)\). The signed power coherence then gives
\(\mathsf{loop}^{\mathsf{pred}(n)}\cdot\mathsf{loop}
 =\mathsf{loop}^n=d(n)\).
Compose these paths at every \(n\), and apply funext. The result has
exactly the type \(T_D(\mathsf{loop},d)=d\), the second input required
by the dependent circle eliminator.
Apply that eliminator with \(d\) and this path to define \(\mathrm{decode}\).
Its point computation is judgmental. Its dependent loop computation
identifies \(\operatorname{apd}_{\mathrm{decode}}(\mathsf{loop})\)
with the supplied lift propositionally. No property of \(\mathrm{decode}\)
was assumed while constructing that lift, and no setness of
\(\mathsf{base}=x\) was used.
\end{proof}
\end{annexitem}

\begin{annexitem}{annex:p07-circle-winding}{thm:p07-circle-winding}{Positive, negative and zero transport calculations}
\begin{proof}
First calculate on the canonical tags of an integer. At zero,
\(\mathsf{loop}^{0}\equiv_{\mathrm j}\mathsf{refl}\), so transport
returns \(m\), which equals \(m+0\).
At \(\mathsf{pos}(0)\), the power is
\(\mathsf{refl}\cdot\mathsf{loop}\). Transport composition says first
transport along reflexivity and then along the loop. Hence it is
\(\mathsf{succ}(m)=m+E(\mathsf{pos}(0))\), by the successor recurrence
and right zero. Suppose the positive tag \(k\) has transport
\(m+E(\mathsf{pos}(k))\). The defining next power appends
\(\mathsf{loop}\), so its transport is
\(\mathsf{succ}(m+E(\mathsf{pos}(k)))\).
The successor recurrence identifies this with
\(m+E(\mathsf{pos}(\mathsf Sk))\). This proves every positive case
by natural induction.

At \(\mathsf{neg}(0)\), the power is
\(\mathsf{refl}\cdot\mathsf{loop}^{-1}\), so inverse monodromy
gives \(\mathsf{pred}(m)=m+E(\mathsf{neg}(0))\).
Suppose the negative tag \(k\) has transport
\(m+E(\mathsf{neg}(k))\). Appending \(\mathsf{loop}^{-1}\)
gives \(\mathsf{pred}(m+E(\mathsf{neg}(k)))\).
The predecessor recurrence identifies this with
\(m+E(\mathsf{neg}(\mathsf Sk))\). This proves every negative case.
The transport-composition comparisons used here are propositional:
\(T(p\cdot q,m)=T(q,T(p,m))\).
For an arbitrary \(n:\mathbb Z\), transport this proved equation
family from \(E(Gn)\) to \(n\) along \(\nu_n\).
Taking \(m=0\) and using \(0+n=n\) proves the winding equation.
Neither induction is induction over unknown circle loops.
\end{proof}
\end{annexitem}

\begin{annexitem}{annex:p07-circle-decode-encode}{thm:p07-circle-decode-encode}{Based identity elimination for the first inverse}
\begin{proof}
Use the motive, dependent on both the endpoint and its path,
\[
Q(x,p)=\bigl(\mathrm{decode}_x
             (T_{\mathrm{code}}(p,0))=p\bigr),
\qquad x:S^1,\quad p:\mathsf{base}=x.
\]
At \(x=\mathsf{base}\), \(p=\mathsf{refl}\), transport gives
\(0\) judgmentally, the decoder's point rule gives
\(\mathsf{loop}^0\) judgmentally, and the zero power gives
\(\mathsf{refl}\) judgmentally. Reflexivity is therefore a term of
this base instance of \(Q\).
Based \(J\) produces a term of \(Q(x,p)\) for every variable
endpoint \(x\) and path \(p\). This is precisely \(\eta_x(p)\).
Applying \(J\) only to paths in the fixed type
\(\mathsf{base}=\mathsf{base}\) would not have this motive and is
not the argument used.
\end{proof}
\end{annexitem}

\begin{annexitem}{annex:p07-circle-encode-decode}{thm:p07-circle-encode-decode}{Circle induction and the missing higher coherence}
\begin{proof}
Use the family
\[
M(x)=\prod_{c:\mathrm{code}(x)}
       \bigl(\mathrm{encode}_x(\mathrm{decode}_x(c))=c\bigr).
\]
At the base, the decoder's judgmental point computation turns this
into
\(\prod_{n:\mathbb Z}
(\mathrm{encode}_{\mathsf{base}}(\mathsf{loop}^n)=n)\).
The winding lemma supplies a term \(H:M(\mathsf{base})\), including
all its signed cases.

We now actually supply the circle loop lift. Let
\(u=T_M(\mathsf{loop},H)\) and \(v=H\).
Both are in \(M(\mathsf{base})\), by the endpoints of the loop.
For every \(n:\mathbb Z\), \(u(n)\) and \(v(n)\) are parallel
identity proofs in the set \(\mathbb Z\), with endpoints
\(\mathrm{encode}_{\mathsf{base}}(\mathrm{decode}_{\mathsf{base}}(n))\)
and \(n\). The set witness proved before introducing the circle
therefore gives a path \(u(n)=v(n)\). Funext applied to these paths
gives \(u=v\), which is exactly \(T_M(\mathsf{loop},H)=H\).
This verifies the required transported higher path; it does not
assert that transport in \(M\) is judgmentally trivial.
The dependent circle eliminator with \(H\) and this lift yields a
section of \(M\), namely \(\epsilon\).
The argument uses setness of the integer fiber, not setness of the
circle's loop space. Hence it is not circular.
\end{proof}
\end{annexitem}

\begin{annexitem}{annex:p07-circle-loop-equivalence}{thm:p07-circle-loop-equivalence}{Equivalence, setness and integer addition}
\begin{proof}
For each \(x\), the maps \(\mathrm{encode}_x\) and
\(\mathrm{decode}_x\) have homotopies
\(\mathrm{decode}_x\mathrm{encode}_x\sim\mathrm{id}\) and
\(\mathrm{encode}_x\mathrm{decode}_x\sim\mathrm{id}\)
by the two preceding proofs.
Theorem~\ref{thm:p07-hott-adjointification} turns this particular
quasi-inverse into the defined half-adjoint equivalence, including its
triangle correction. Thus the displayed equivalence is in the actual
equivalence interface, not an undeclared use of a quasi-inverse as a
definition of equivalence.
At the base the code fiber is judgmentally \(\mathbb Z\).
Equivalence preserves setness by
Theorem~\ref{thm:p07-hott-level-closure}; therefore the loop space
is a set. Encoding reflexivity is judgmentally zero; decoding zero
is judgmentally reflexivity.

For \(p,q:\mathsf{base}=\mathsf{base}\), transport composition gives
\[
\mathrm{encode}(p\cdot q)
 =T_{\mathrm{code}}(q,\mathrm{encode}(p)).
\]
The inverse of \(\eta_{\mathsf{base}}(q)\), together with the decoder
point computation, identifies \(q\) with
\(\mathsf{loop}^{\mathrm{encode}(q)}\).
Apply transport to this path equality. The winding lemma, now with
\(m=\mathrm{encode}(p)\), gives
\[
T_{\mathrm{code}}(q,\mathrm{encode}(p))
 =\mathrm{encode}(p)+\mathrm{encode}(q).
\]
This proves the claimed operation comparison. It also transfers the
loop unit, association and inverse laws to integer addition, if needed;
no independent group-structure axioms about the quotient were assumed.
\end{proof}
\end{annexitem}

\begin{annexitem}{annex:p07-circle-cover-contractible}{thm:p07-circle-cover-contractible}{A contraction after encode--decode}
\begin{proof}
Choose center \((\mathsf{base},0)\).
For \((x,c):\sum_x\mathrm{code}(x)\), the Sigma-path constructor of
Theorem~\ref{thm:p07-hott-sigma-paths} takes the first coordinate
\(p=\mathrm{decode}_x(c):\mathsf{base}=x\) and the second coordinate
\[
\epsilon_x(c):
T_{\mathrm{code}}(\mathrm{decode}_x(c),0)=c.
\]
The latter has exactly the required type, by the definition of encoding.
Their paired path is a path from the center to \((x,c)\).
This is a dependent contraction and proves contractibility.
It was derived after constructing both inverse homotopies; it is not
being used to justify either one, nor does it prove the general
flattening lemma used in the source's alternative route.
\end{proof}
\end{annexitem}

\begin{annexitem}{annex:p07-circle-power-coherence}{thm:p07-circle-power-coherence}{All signed power boundary cases}
\begin{proof}
It suffices to work on \(E(c)\): afterward transport the equation family
along \(\nu_n:E(Gn)=n\).
For the successor equation at zero the displayed definition gives
\(\ell^{1}=\mathsf{refl}\cdot\ell\).
At a positive tag the defining positive recursion gives it.
At \(\mathsf{neg}(0)\), the required equation is
\(\mathsf{refl}=(\mathsf{refl}\cdot\ell^{-1})\cdot\ell\).
Associate, use \(\ell^{-1}\cdot\ell=\mathsf{refl}\), and apply the
unit laws. At \(\mathsf{neg}(\mathsf Sk)\), the required equation is
\[
P_\ell(\mathsf{neg}(k))
 =(P_\ell(\mathsf{neg}(k))\cdot\ell^{-1})\cdot\ell,
\]
proved by the same association, inverse cancellation and right-unit
paths. These are all successor cases.

For the predecessor equation at zero use the negative initial clause;
at any negative tag use the negative recursion. At \(\mathsf{pos}(0)\)
it reads
\(\mathsf{refl}=(\mathsf{refl}\cdot\ell)\cdot\ell^{-1}\),
and at \(\mathsf{pos}(\mathsf Sk)\) it reads
\[
P_\ell(\mathsf{pos}(k))
 =(P_\ell(\mathsf{pos}(k))\cdot\ell)\cdot\ell^{-1}.
\]
Both follow by association, \(\ell\cdot\ell^{-1}=\mathsf{refl}\),
and right unit. These are all predecessor cases.

For the decoder's coherence directly: at zero its left side is
\((\mathsf{refl}\cdot\ell^{-1})\cdot\ell\); at
\(\mathsf{pos}(0)\) it is \(\mathsf{refl}\cdot\ell\); at
\(\mathsf{pos}(\mathsf Sk)\) it is
\(P_\ell(\mathsf{pos}(k))\cdot\ell\).
The first reduces propositionally by inverse cancellation and units;
the next two are the power's defining positive clauses.
At \(\mathsf{neg}(k)\) it is
\((P_\ell(\mathsf{neg}(k))\cdot\ell^{-1})\cdot\ell\),
which equals \(P_\ell(\mathsf{neg}(k))\) by association, inverse
cancellation and right unit. All invoked paths are instances of
Theorem~\ref{thm:p07-hott-path-laws}. None uses setness of the loop
space, which has not yet been proved.
\end{proof}
\end{annexitem}

\begin{annexitem}{annex:p07-circle-monodromy}{thm:p07-circle-monodromy}{Univalent monodromy and its inverse}
\begin{proof}
The pullback transport comparison of
Theorem~\ref{thm:p07-hott-transport-laws} identifies
\(T_{\mathrm{code}}(\mathsf{loop},n)\) with
\(T_{X\mapsto X}(\operatorname{ap}_{\mathrm{code}}(\mathsf{loop}),n)\).
Apply the propositional loop computation of circle recursion to replace
the universe path by \(\operatorname{ua}(e_{\mathsf{succ}})\).
The univalence comparison of
Theorem~\ref{thm:p07-hott-univalence-comparison} identifies this last
transport with the forward map of \(e_{\mathsf{succ}}\), namely
\(\mathsf{succ}(n)\). This proves forward monodromy, with no
judgmental univalence assertion.

For the inverse calculation put
\(u=T_{\mathrm{code}}(\mathsf{loop}^{-1},n)\).
Inverse transport gives
\(T_{\mathrm{code}}(\mathsf{loop},u)=n\).
Forward monodromy yields \(\mathsf{succ}(u)=n\).
Apply \(\operatorname{ap}_{\mathsf{pred}}\), then compose with the
inverse of the previously proved path
\(\mathsf{pred}(\mathsf{succ}(u))=u\), to obtain
\(u=\mathsf{pred}(n)\).
This proves the particular inverse transport needed by the decoder,
rather than relying on an unspecified inverse of a family.
\end{proof}
\end{annexitem}
